% !PRINTED 394
\weberpage{394}
%<<CONT
\[
\frac{5 + \sqrt{37}}{2} = [5, 1, 1].
\]
5. $D = 148 = 4 \ldotp 37 \quad \lambda = 12$
\[
b = \quad 2, \quad 4, \quad 6, \quad 8, \quad 10, \quad 12
\]
\[
a\,c = 36, \quad 33, \quad 28, \quad 21, \quad 12, \quad 1;
\]
\gesperrt{limits} for $a$, $c$:
\[
6, 7; \quad 5, 8; \quad 4, 9; \quad 3, 10; \quad 2, 11; \quad 1, 12;
\]
\gesperrt{reduced numbers}:
\[
\{4, 6, 7\}, \quad \{7, 6, 4\}, \quad \{3, 8, 7\}, \quad \{7, 8, 3\}
\]
\[
\{3, 10, 4\}, \quad \{4, 10, 3\}, \quad \{1, 12, 1\}.
\]
Here we obtain three periods of 1, 3, and 3 terms:
\[
\{1, 12, 1\}, \quad \{1, 12, 1\},
\]
\[
\{3, 10, 4\}, \quad \{4, 6, 7\}, \quad \{7, 8, 3\}, \quad \{3, 10, 4\},
\]
\[
\{4, 10, 3\}, \quad \{3, 8, 7\}, \quad \{7, 6, 4\}, \quad \{4, 10, 3\}
\]
and the continued-fraction periods
\[
[12], \quad [2, 1, 3], \quad [3, 1, 2];
\]
of these the first is reversible, and the other two pass into one
another under reversal.

6. $D = 136 = 4 \ldotp 34, \quad \lambda = 11,$
\[
b = \quad 2, \quad 4, \quad 6, \quad 8, \quad 10,
\]
\[
a\,c = 33, \quad 30, \quad 25, \quad 18, \quad 9;
\]
\gesperrt{limits} for $a$, $c$:
\[
5, 6; \quad 4, 7; \quad 3, 8; \quad 2, 9; \quad 1, 10;
\]
\gesperrt{reduced numbers}:
\[
\{5, 4, 6\}, \quad \{6, 4, 5\}, \quad \{5, 6, 5\}, \quad \{2, 8, 9\}, \quad \{9, 8, 2\}
\]
\[
\{3, 8, 6\}, \quad \{6, 8, 3\}, \quad \{1, 10, 9\}, \quad \{9, 10, 1\}, \quad \{3, 10, 3\};
\]
\gesperrt{periods}:
\[
\{5, 4, 6\}, \quad \{6, 8, 3\}, \quad \{3, 10, 3\}, \quad \{3, 8, 6\}, \quad \{6, 4, 5\}, \quad \{5, 6, 5\}, \quad \{5, 4, 6\},
\]
\[
\{9, 10, 1\}, \quad \{1, 10, 9\}, \quad \{9, 8, 2\}, \quad \{2, 8, 9\}, \quad \{9, 10, 1\};
\]
\gesperrt{continued-fraction periods}:
\[
[1, 3, 3, 1, 1, 1], \quad [10, 1, 4, 1].
\]
Here, then, we have two different periods, both of which are
reversible.
